\documentclass[10pt,a4paper]{amsart}
\usepackage[margin=19mm]{geometry}
\usepackage{amsmath,amssymb,mathtools}
\usepackage[scr=rsfs,cal=euler]{mathalfa}
\usepackage{xfrac}
\usepackage[unicode,hidelinks,bookmarks=false]{hyperref}
\newcommand{\ie}{i.e.\ }
\newcommand{\cf}{cf.\ }
\newcommand{\resp}{respectively\ }
\newcommand{\Subsection}{\subsection}
\providecommand{\nobreakdash}{\nobreak\hbox{-}}
\setlength{\emergencystretch}{2em}
\multlinegap0pt
\usepackage{bm}
\usepackage{bbm}
\usepackage{dsfont}
\usepackage{xspace}
\usepackage{cancel}
\usepackage{mathtools}
\usepackage{amsthm}
\makeatletter
\@ifundefined{proposition}{\newtheorem{proposition}{Proposition}}{}
\makeatother
\newcommand{\Aspec}{A_{\text{spec}}}
\newcommand{\R}{\mathbb{R}}
\newcommand{\bvec}{\mathbf{b}}
\newcommand{\xbase}{\mathbf{x}_{\text{base}}}
\newcommand{\xspec}{\mathbf{x}_{\text{spec}}}
\title[Contextual Deconvolution for Variance-Stable Demand Sensing:]{Contextual Deconvolution for Variance-Stable Demand Sensing: Kernel-Modulated Operators in Promotional Retail}
\author{Mohammad Forouhesh}
\date{Source: arXiv:2607.25664v1 (CC BY 4.0)}
\begin{document}
\maketitle

\noindent\emph{Source and licence.} Contextual Deconvolution for Variance-Stable Demand Sensing: Kernel-Modulated Operators in Promotional Retail. Version arXiv:2607.25664v1 (CC BY 4.0), distributed under the Creative Commons Attribution 4.0 International licence, as recorded in the arXiv licence metadata for this exact version and shown on the abstract page https://arxiv.org/abs/2607.25664v1. Full rights notice: licenses/arxiv-2607.25664-CC-BY-4.0.txt.\par\smallskip

\noindent\textit{Editorial scope.} Counted unit: the Proposition whose source label is \texttt{None} in the article (source0.tex lines 523--532, source label \texttt{None}), delivered together with the article's own complete original proof of it (source0.tex lines 534--562, one \texttt{proof} environment of 29 lines). No mathematics is written, repaired or re-derived: statement and proof are the article's own text line for line. Required context retained uncounted: none. Every definition, display and label of those lines is retained unchanged. Omitted from this package and not redistributed: 1--522, 533--533, 563--1981. Nothing inside a retained excerpt is omitted or altered: every retained body is delivered exactly as the article's lines are written, so no omission receipt of any kind accompanies this package. Citations: the retained excerpts carry no citation command at all, so no bibliography entry of the article is transcribed and no reference is redistributed. Reference adaptation: none was needed and none was made; every reference inside the delivered unit resolves inside it, so no unresolved reference or citation is produced. Printed slips kept literal: no printed word, symbol, hypothesis or number of the counted statement or of its proof was altered. Layout and class: the article's own class is \texttt{article} with packages including tmlr, amsmath, amssymb, amsthm, mathtools, bm, graphicx, booktabs, array, multirow, enumitem, hyperref, url; the wrapper is the lane's pinned \texttt{amsart} editorial wrapper at 10pt on A4, which restates the theorem-like environments the retained text needs and adds the article's own macro definitions that the retained text needs (\texttt{\textbackslash Aspec}, \texttt{\textbackslash R}, \texttt{\textbackslash bvec}, \texttt{\textbackslash xbase}, \texttt{\textbackslash xspec}), with extra wrapper packages (bm, bbm, dsfont, xspace, cancel, mathtools). The delivered numbering is the unit's own. Assets: no retained span contains a figure environment, a graphics-inclusion command, a PDF-page inclusion, a table or a TikZ picture; no image file is copied into this package, the package declares no asset dependency and no image file is redistributed with it. Licence: version arXiv:2607.25664v1 of this article is distributed under the Creative Commons Attribution 4.0 International licence, as recorded in the arXiv licence metadata for this exact version and shown on the abstract page https://arxiv.org/abs/2607.25664v1. A full rights notice is delivered at licenses/arxiv-2607.25664-CC-BY-4.0.txt. Characters: every retained excerpt is pure ASCII, so the delivered engine, XeLaTeX with its default Latin Modern text font, reports no missing character.
\begin{proposition}[Convexity]
For any fixed, causal convolution operator $\Aspec \in \R^{T \times T}$ with finite bandwidth $n$, the objective
\begin{equation}
\begin{split}
J(\xbase, \xspec) &= \|\bvec - (\xbase + \Aspec \xspec)\|_2^2 \\
&\quad + \lambda_{\text{smooth}} \|L \xbase\|_2^2 + \lambda_{\text{sparse}} \|\xspec\|_1
\end{split}
\end{equation}
is convex in $(\xbase, \xspec)$.
\end{proposition}

\begin{proof}
Decompose $J = f + g$ where
\begin{align*}
f(\xbase, \xspec) &= \|\bvec - \xbase - \Aspec \xspec\|_2^2 + \lambda_{\text{smooth}} \|L \xbase\|_2^2, \\
g(\xspec) &= \lambda_{\text{sparse}} \|\xspec\|_1.
\end{align*}

\textbf{Step 1: $f$ is convex.} The squared Euclidean norm $\|\cdot\|_2^2$ is convex. Composing it with an affine map preserves convexity. The map $(\xbase, \xspec) \mapsto \bvec - \xbase - \Aspec \xspec$ is affine in $(\xbase, \xspec)$. Therefore $\|\bvec - \xbase - \Aspec \xspec\|_2^2$ is convex. Similarly, $\|L \xbase\|_2^2$ is convex because $L$ is linear and the squared norm is convex. The sum of convex functions is convex, so $f$ is convex.

For completeness, the Hessian of $f$ with respect to $(\xbase, \xspec)$ is the symmetric block matrix
\begin{equation}
H = \begin{bmatrix} I + \lambda_{\text{smooth}} L^\top L & \Aspec \\ \Aspec^\top & \Aspec^\top \Aspec \end{bmatrix}.
\end{equation}

To verify $H \succeq 0$ directly, note that the upper-left block $A = I + \lambda_{\text{smooth}} L^\top L$ is positive definite. The Schur complement of $A$ in $H$ is
\begin{equation}
S = \Aspec^\top \Aspec - \Aspec^\top (I + \lambda_{\text{smooth}} L^\top L)^{-1} \Aspec.
\end{equation}

Since $A = I + \lambda_{\text{smooth}} L^\top L \succeq I$, we have $A^{-1} \preceq I$, which implies
\begin{equation}
\Aspec^\top A^{-1} \Aspec \preceq \Aspec^\top \Aspec.
\end{equation}
Thus $S \succeq 0$. Because $A \succ 0$ and $S \succeq 0$, the block matrix $H$ is positive semidefinite.

\textbf{Step 2: $g$ is convex.} The $\ell_1$ norm is convex on $\mathbb{R}^T$. Scaling by $\lambda_{\text{sparse}} > 0$ preserves convexity.

\textbf{Step 3: Sum of convex functions.} The sum of convex functions is convex. Therefore $J = f + g$ is convex in $(\xbase, \xspec)$.
\end{proof}

\end{document}
